
% How does our experience and results differ from prior work? I think one distinguishing property is the decentralized nature of build infrastructure, introducing lots of diversity that make recovering rebuild metadata difficult. 

% Two questions: Can we recover the data we need if it is not recorded, and can we verify the artifact was built from the source? The major limitation will be in recovering the build instructions and dealing with monorepos with complex build steps

% I think the core question here is: How should we define L3/L4? Should we assume safety if no new files are added?

% When should verification happen? There are failures that happened because a dependency version was no longer available. Should verification be done at package upload time and a badge presented to users? Should it be done by the users? The first reduces dependency issues. The second gives users more control power.

% Maybe we can frame the question "Should developers deviate from default practices?" Should there be some unifying strategies? For example, developers can declare different build commands and keys but we need build info file that produces the sequences of commands that lead to a reproducible build? 

% I think this paper is asking the question: "Is it possible to recover the build metadata for the original build such that a rebuilt artifact is equivalent?

% Papers on attesting pipelines
% Attesting LLM Pipelines: Enforcing Verifiable Training and Release Claims
% Kettle: Attested Builds for Verifiable Software Provenance
% AEX: Non-Intrusive Multi-Hop Attestation and Provenance for LLM APIs

% Standards focus on making builds reproducible. Is this enough?

\section{Discussion}

Open-source software supply-chain standards recommend reproducible builds as a means of improving trust in distributed software artifacts.
We measured the gap between these recommendations and artifact verifiability in practice.

\subsection{Verifiability Requires Artifact-Bound Rebuild Metadata}
\label{subsec:discussion-verifiability-metadata}

\Cref{subsec:results-rq1} shows that end-to-end verification (over all sampled artifacts) ranges between 25\% and 80\% across ecosystems.
\Cref{subsec:results-rq2} further shows that many failures stem from unpinned build tools, CI/CD source modifications, and custom packaging.
Independent verifiability therefore requires more than deterministic builds: it needs durable, structured, artifact-bound or recoverable metadata that lets independent parties rebuild and compare the published artifact.

Centralized build ecosystems such as Debian, Arch Linux, and Nix partly address this problem by recording rebuild metadata through ecosystem-controlled build infrastructure and providing verifiers that consume them.
This model is difficult to apply directly to decentralized-build ecosystems where maintainers use diverse build tools, packaging conventions, dependency managers, and CI/CD workflows.
A more practical direction is for other centralized infrastructure, such as CI/CD workflow runners, to record rebuild metadata while executing workflow steps and embed in artifacts before publishing.
Our replay pilot (\Cref{subsec:results-rq2}) reinforces this: the one artifact recovered by re-running its workflow matched only because the workflow pinned the build-toolchain version our heuristic missed---it is the build environment, not the workflow logic, that must be recorded.
Additionally, developer practices that improve verifiability, such as pinning dependencies and including package preparation and build steps in package manifest or other metadata files, should be encouraged.

This also bears on supply-chain standards. Existing standards recommend reproducible builds~\cite{noauthor_securing_nodate}, but our results show reproducibility becomes actionable only when ecosystems also make recovering rebuild metadata feasible; standards should therefore add concrete guidance on which metadata to preserve, where verifiers obtain it, and how results are represented, informed by our RQ2 causes (\cref{subsec:results-rq2}).

\subsection{Deploying Verifiability in Practice}
\label{subsec:discussion-deploying-verifiability}

This paper reframes reproducible builds as a necessary but incomplete condition for software supply-chain integrity in decentralized-build ecosystems.
The practical barrier is whether an independent verifier can recover enough rebuild metadata to reproduce the published artifact. 
By identifying the metadata gaps and packaging practices that prevent verification, this study gives registries, standards bodies, and rebuild-tool developers concrete targets for making artifact verification practical at package-registry scale.

One challenge is how verifiability can be integrated into package lifecycles without immediate, ecosystem-wide enforcement. 
Registries are a natural deployment point---they already receive, validate, and distribute artifacts---and could adopt it incrementally: first publishing verification metadata and warnings, then prioritizing stronger equivalence levels for popular, critical, or security-sensitive packages. Downstream package managers and organizational policies can require a target equivalence level (\cref{subsec:artifact-comparison-model}) before installation to protect users.

Provenance attestations offer another path: they improve metadata recovery but attest only to the build process, not the artifact's authenticity~\cite{openssf_introducing_nodate}. They could instead distinguish \emph{workflow} provenance from \emph{package-build} provenance providing sufficient evidence to reproduce the artifact~\cite{noauthor_build_nodate}.


% \subsection{Directions for Software Engineering Research}
% \label{subsec:discussion-research-directions}

% Our findings raise three directions for software engineering research.
% First, researchers should study verification-aware attestations that record artifact-level rebuild metadata while accounting for ecosystem diversity, maintainer burden, and trusted publishing overhead.
% This includes identifying which source, dependency, environment, and build-command metadata should be recorded and how attestations should represent verification results.
% Second, more work is needed to understand why provenance and trusted publishing adoption remain low, even among recent and popular packages, and which ecosystem incentives or usability improvements can increase adoption.
% Third, future work should develop security-aware comparison models that distinguish benign metadata differences from security-relevant changes.
% The tiered model used in this paper is one step in this direction, but additional work is needed to identify which differences are safe to accommodate at L3 and how to operationalize L4 behavioral equivalence without masking malicious changes.


% \section{Discussion}

% Open-source software supply chain standards increasingly recommend reproducible builds and artifact verifiability as mechanisms for improving trust in distributed software artifacts. 
% Our results expose a gap between these recommendations and the infrastructure available in source-oriented package ecosystems.
% We discuss the implications of this gap for software engineering practice, supply-chain security, and future research.

% \subsection{Addressing Ecosystem-Level Metadata Gaps for Verifiability}
% \label{subsec:discussion-verifiability-metadata}

% \Cref{subsec:results-rq1} shows that, although reproducibility rates exceed XX\% when the necessary rebuild metadata is available, artifact verifiability rates remain low across ecosystems.
% \Cref{subsec:results-rq2} further shows that many failures stem from missing source metadata, custom build practices, and environment-dependent build pipelines that are difficult to reconstruct automatically.
% These results suggest that independent verifiability requires more than deterministic builds: it requires durable, structured, and artifact-bound metadata that allows an external verifier to recover the source, build environment, and packaging process.

% Centralized build ecosystems such as Debian, Arch Linux, and Nix address this problem by modifying their build infrastructure to record rebuild metadata and bind it to produced artifacts.
% This model is difficult to apply directly to source-oriented ecosystems such as npm and PyPI, where packages use diverse build tools, packaging conventions, dependencies, and CI/CD workflows.
% A more practical direction is to capture rebuild metadata at the CI/CD layer.
% For example, CI/CD runners such as GitHub Actions could record the workflow configuration, executed build steps, source state, environment, and artifact-producing commands, then bind this information to the published artifact.
% Such direction would address the generalizability limits of using canonical build commands and the security risks of replaying workflow steps discussed in \cref{subsec:verifier-model}, make verifiability more practical for packages that already build and publish through CI/CD and encourage other verifiability-interested maintainers to adopt CI/CD and trusted publishing pipelines.

% \subsection{Artifact Verifiability in Software Supply Chain Standards}

% Software supply-chain frameworks and standards increasingly identify reproducible builds as a mechanism for detecting compromised artifacts.
% These recommendations have motivated research and industry efforts to improve reproducibility across software ecosystems, and several ecosystems and organizations now report reproducible artifacts.
% However, our results show that, in source-oriented package ecosystems, higher reproducibility does not automatically translate into higher verifiability.
% Reproducibility is necessary for artifact verifiability, but it becomes a practical supply-chain defense only when parties other than the original builder can independently recover the required rebuild metadata, reproduce the artifact, and compare it against the distributed package.

% This gap suggests that existing standards should complement reproducibility recommendations with more concrete guidance for independent artifact verification.
% In particular, standards should specify what rebuild metadata ecosystems should preserve, where that metadata should be exposed, how it should be bound to published artifacts, and how verifiers should consume it.
% Our results provide empirical evidence for these requirements by identifying where verifiability fails and which missing metadata, build practices, and ecosystem characteristics contribute to those failures.


% \subsection{Enhancing Provenance Attestations for Verification}
% \label{subsec:discussion-attestation-verification}

% Our results show that provenance attestations improve verifiability by providing higher-fidelity source and workflow information than ad hoc registry metadata.
% However, they do not, by themselves, make artifacts independently verifiable.
% Current attestations often identify the workflow trigger commit and CI/CD workflow, but this information may not capture the exact source tree or source commit, build configuration, or artifact-producing commands used to build the package.
% In addition, workflow URLs may be unavailable, mutable, or insufficient for reconstructing the build process.
% As a result, current attestations can provide evidence that an artifact was produced in a particular CI/CD context, but they do not necessarily prove that the artifact corresponds to the intended source or that the build process was uncompromised.
% This limitation is especially important given recent supply-chain incidents involving compromised CI/CD workflows and GitHub Actions.

% Improving trusted publishing for artifact verification requires distinguishing workflow provenance from package-build provenance.
% Workflow provenance records the CI/CD context in which a build occurred, including the runner, workflow, and trigger event.
% Package-build provenance should record the exact source tree, build configuration, dependency state, and commands used to produce the published package.
% If the artifact is independently verified during publishing, the attestation could also include the verification result and equivalence level.
% This design would make attestations more directly useful for independent verification, reduce dependence on ephemeral source-hosting state, and provide more durable evidence for registries, downstream users, and third-party auditing services.


% \subsection{Deploying Verifiability in Practice}
% \label{subsec:discussion-deploying-verifiability}

% Artifact verifiability can strengthen defenses against attacks that compromise build, release, or publication infrastructure.
% As shown in \cref{sec:knowledge-gaps}, many recent incidents introduced malicious code after, or outside, normal source development workflows.
% In these cases, verifiability provides a way to check whether the distributed artifact matches the expected source.
% The practical challenge is how to integrate this check into existing package lifecycles.

% Registries are a natural deployment point because they already receive, validate, and distribute artifacts.
% Registry-side verification could reject suspicious artifacts, label unverifiable artifacts, or expose verification status to downstream users.
% However, our results show that immediate hard enforcement would be impractical because many packages are not yet independently verifiable.
% A more realistic path is incremental deployment: registries could first publish verification metadata and warnings, then prioritize stronger requirements for popular, critical, or security-sensitive packages.

% Downstream users and organizations are another deployment point because they bear the risk of compromised dependencies.
% Package managers and organizational dependency policies could allow users to require artifacts verified at specific equivalence levels.
% The tiered comparison model in \cref{subsec:artifact-comparison-model} supports this deployment model by avoiding an all-or-nothing definition of reproducibility.
% Some environments may require byte-identical artifacts, while others may accept artifacts that differ only in benign metadata files or fields.
% This allows verification policies to reflect different risk tolerances while still encouraging maintainers to improve their build and release practices.

% These deployment paths also suggest a role for supply-chain frameworks such as SLSA.
% A verified-build profile could complement SLSA level 3's hardened build systems by requiring evidence that artifacts can be independently rebuilt and compared by an alternative verifier.
% Such a profile would better address attacks in which the original build or release infrastructure produces a compromised artifact despite satisfying provenance, isolation, or trusted-publishing requirements.


% \subsection{Directions for Software Engineering Research}
% \label{subsec:discussion-research-directions}

% Our findings suggest three directions for research:

% \myparagraph{1. Verification-aware attestations} 
% Because current attestations improve source recovery but do not consistently provide artifact-level rebuild metadata, researchers should study how provenance attestations can be redesigned for artifact verification.
% This includes identifying which source, dependency, environment, and build-command metadata should be recorded; how much overhead verification adds to trusted publishing; and how attestations should represent verification results and equivalence levels.
% Recent work on verifying software systems using in-toto attestations provides useful starting points, but artifact verification in source-oriented package ecosystems must also account for diverse build tools, packaging conventions, and release practices.

% \myparagraph{2. Adoption of trusted publishing and provenance}
% Our results show low provenance-attestation adoption across ecosystems, even among recent and popular packages.
% More work is needed to understand which factors affect adoption and which ecosystem strategies can improve it.
% Prior work has studied signing, build-pipeline trust, and provenance in specific systems, but less is known about which ecosystem features encourage adoption, which maintainer costs discourage it, and whether trusted publishing changes build and release practices over time.
% Answering these questions is important because provenance can improve verifiability only if it is widely adopted and consistently bound to published artifacts.

% \myparagraph{3. Security-aware comparison models}
% Finally, our root-cause analysis shows that artifact verifiability is affected by diverse developer practices.
% Much prior work on reproducible builds recommends changes to developer practices or infrastructure to eliminate sources of non-reproducibility.
% However, such changes can be slow to adopt because maintainers face competing priorities, workflow constraints, and organizational inertia.
% This raises an important research question: which sources of artifact differences should be eliminated through better infrastructure, and which benign differences should be accommodated by comparison models?
% The tiered model in \cref{subsec:artifact-comparison-model} is one step in this direction because it permits lower equivalence levels for artifacts that differ only in known metadata files or fields.
% Future work should develop principled methods for identifying which differences are safe to accommodate at L3 and techniques for operationalizing L4 behavioral equivalence without masking security-relevant changes.
